% Options for packages loaded elsewhere
\PassOptionsToPackage{unicode}{hyperref}
\PassOptionsToPackage{hyphens}{url}
\PassOptionsToPackage{dvipsnames,svgnames,x11names}{xcolor}
%
\documentclass[
  letterpaper,
  DIV=11,
  numbers=noendperiod]{scrartcl}

\usepackage{amsmath,amssymb}
\usepackage{iftex}
\ifPDFTeX
  \usepackage[T1]{fontenc}
  \usepackage[utf8]{inputenc}
  \usepackage{textcomp} % provide euro and other symbols
\else % if luatex or xetex
  \usepackage{unicode-math}
  \defaultfontfeatures{Scale=MatchLowercase}
  \defaultfontfeatures[\rmfamily]{Ligatures=TeX,Scale=1}
\fi
\usepackage{lmodern}
\ifPDFTeX\else  
    % xetex/luatex font selection
\fi
% Use upquote if available, for straight quotes in verbatim environments
\IfFileExists{upquote.sty}{\usepackage{upquote}}{}
\IfFileExists{microtype.sty}{% use microtype if available
  \usepackage[]{microtype}
  \UseMicrotypeSet[protrusion]{basicmath} % disable protrusion for tt fonts
}{}
\makeatletter
\@ifundefined{KOMAClassName}{% if non-KOMA class
  \IfFileExists{parskip.sty}{%
    \usepackage{parskip}
  }{% else
    \setlength{\parindent}{0pt}
    \setlength{\parskip}{6pt plus 2pt minus 1pt}}
}{% if KOMA class
  \KOMAoptions{parskip=half}}
\makeatother
\usepackage{xcolor}
\setlength{\emergencystretch}{3em} % prevent overfull lines
\setcounter{secnumdepth}{-\maxdimen} % remove section numbering
% Make \paragraph and \subparagraph free-standing
\makeatletter
\ifx\paragraph\undefined\else
  \let\oldparagraph\paragraph
  \renewcommand{\paragraph}{
    \@ifstar
      \xxxParagraphStar
      \xxxParagraphNoStar
  }
  \newcommand{\xxxParagraphStar}[1]{\oldparagraph*{#1}\mbox{}}
  \newcommand{\xxxParagraphNoStar}[1]{\oldparagraph{#1}\mbox{}}
\fi
\ifx\subparagraph\undefined\else
  \let\oldsubparagraph\subparagraph
  \renewcommand{\subparagraph}{
    \@ifstar
      \xxxSubParagraphStar
      \xxxSubParagraphNoStar
  }
  \newcommand{\xxxSubParagraphStar}[1]{\oldsubparagraph*{#1}\mbox{}}
  \newcommand{\xxxSubParagraphNoStar}[1]{\oldsubparagraph{#1}\mbox{}}
\fi
\makeatother

\usepackage{color}
\usepackage{fancyvrb}
\newcommand{\VerbBar}{|}
\newcommand{\VERB}{\Verb[commandchars=\\\{\}]}
\DefineVerbatimEnvironment{Highlighting}{Verbatim}{commandchars=\\\{\}}
% Add ',fontsize=\small' for more characters per line
\usepackage{framed}
\definecolor{shadecolor}{RGB}{241,243,245}
\newenvironment{Shaded}{\begin{snugshade}}{\end{snugshade}}
\newcommand{\AlertTok}[1]{\textcolor[rgb]{0.68,0.00,0.00}{#1}}
\newcommand{\AnnotationTok}[1]{\textcolor[rgb]{0.37,0.37,0.37}{#1}}
\newcommand{\AttributeTok}[1]{\textcolor[rgb]{0.40,0.45,0.13}{#1}}
\newcommand{\BaseNTok}[1]{\textcolor[rgb]{0.68,0.00,0.00}{#1}}
\newcommand{\BuiltInTok}[1]{\textcolor[rgb]{0.00,0.23,0.31}{#1}}
\newcommand{\CharTok}[1]{\textcolor[rgb]{0.13,0.47,0.30}{#1}}
\newcommand{\CommentTok}[1]{\textcolor[rgb]{0.37,0.37,0.37}{#1}}
\newcommand{\CommentVarTok}[1]{\textcolor[rgb]{0.37,0.37,0.37}{\textit{#1}}}
\newcommand{\ConstantTok}[1]{\textcolor[rgb]{0.56,0.35,0.01}{#1}}
\newcommand{\ControlFlowTok}[1]{\textcolor[rgb]{0.00,0.23,0.31}{\textbf{#1}}}
\newcommand{\DataTypeTok}[1]{\textcolor[rgb]{0.68,0.00,0.00}{#1}}
\newcommand{\DecValTok}[1]{\textcolor[rgb]{0.68,0.00,0.00}{#1}}
\newcommand{\DocumentationTok}[1]{\textcolor[rgb]{0.37,0.37,0.37}{\textit{#1}}}
\newcommand{\ErrorTok}[1]{\textcolor[rgb]{0.68,0.00,0.00}{#1}}
\newcommand{\ExtensionTok}[1]{\textcolor[rgb]{0.00,0.23,0.31}{#1}}
\newcommand{\FloatTok}[1]{\textcolor[rgb]{0.68,0.00,0.00}{#1}}
\newcommand{\FunctionTok}[1]{\textcolor[rgb]{0.28,0.35,0.67}{#1}}
\newcommand{\ImportTok}[1]{\textcolor[rgb]{0.00,0.46,0.62}{#1}}
\newcommand{\InformationTok}[1]{\textcolor[rgb]{0.37,0.37,0.37}{#1}}
\newcommand{\KeywordTok}[1]{\textcolor[rgb]{0.00,0.23,0.31}{\textbf{#1}}}
\newcommand{\NormalTok}[1]{\textcolor[rgb]{0.00,0.23,0.31}{#1}}
\newcommand{\OperatorTok}[1]{\textcolor[rgb]{0.37,0.37,0.37}{#1}}
\newcommand{\OtherTok}[1]{\textcolor[rgb]{0.00,0.23,0.31}{#1}}
\newcommand{\PreprocessorTok}[1]{\textcolor[rgb]{0.68,0.00,0.00}{#1}}
\newcommand{\RegionMarkerTok}[1]{\textcolor[rgb]{0.00,0.23,0.31}{#1}}
\newcommand{\SpecialCharTok}[1]{\textcolor[rgb]{0.37,0.37,0.37}{#1}}
\newcommand{\SpecialStringTok}[1]{\textcolor[rgb]{0.13,0.47,0.30}{#1}}
\newcommand{\StringTok}[1]{\textcolor[rgb]{0.13,0.47,0.30}{#1}}
\newcommand{\VariableTok}[1]{\textcolor[rgb]{0.07,0.07,0.07}{#1}}
\newcommand{\VerbatimStringTok}[1]{\textcolor[rgb]{0.13,0.47,0.30}{#1}}
\newcommand{\WarningTok}[1]{\textcolor[rgb]{0.37,0.37,0.37}{\textit{#1}}}

\providecommand{\tightlist}{%
  \setlength{\itemsep}{0pt}\setlength{\parskip}{0pt}}\usepackage{longtable,booktabs,array}
\usepackage{calc} % for calculating minipage widths
% Correct order of tables after \paragraph or \subparagraph
\usepackage{etoolbox}
\makeatletter
\patchcmd\longtable{\par}{\if@noskipsec\mbox{}\fi\par}{}{}
\makeatother
% Allow footnotes in longtable head/foot
\IfFileExists{footnotehyper.sty}{\usepackage{footnotehyper}}{\usepackage{footnote}}
\makesavenoteenv{longtable}
\usepackage{graphicx}
\makeatletter
\def\maxwidth{\ifdim\Gin@nat@width>\linewidth\linewidth\else\Gin@nat@width\fi}
\def\maxheight{\ifdim\Gin@nat@height>\textheight\textheight\else\Gin@nat@height\fi}
\makeatother
% Scale images if necessary, so that they will not overflow the page
% margins by default, and it is still possible to overwrite the defaults
% using explicit options in \includegraphics[width, height, ...]{}
\setkeys{Gin}{width=\maxwidth,height=\maxheight,keepaspectratio}
% Set default figure placement to htbp
\makeatletter
\def\fps@figure{htbp}
\makeatother

\KOMAoption{captions}{tableheading}
\makeatletter
\@ifpackageloaded{caption}{}{\usepackage{caption}}
\AtBeginDocument{%
\ifdefined\contentsname
  \renewcommand*\contentsname{Table of contents}
\else
  \newcommand\contentsname{Table of contents}
\fi
\ifdefined\listfigurename
  \renewcommand*\listfigurename{List of Figures}
\else
  \newcommand\listfigurename{List of Figures}
\fi
\ifdefined\listtablename
  \renewcommand*\listtablename{List of Tables}
\else
  \newcommand\listtablename{List of Tables}
\fi
\ifdefined\figurename
  \renewcommand*\figurename{Figure}
\else
  \newcommand\figurename{Figure}
\fi
\ifdefined\tablename
  \renewcommand*\tablename{Table}
\else
  \newcommand\tablename{Table}
\fi
}
\@ifpackageloaded{float}{}{\usepackage{float}}
\floatstyle{ruled}
\@ifundefined{c@chapter}{\newfloat{codelisting}{h}{lop}}{\newfloat{codelisting}{h}{lop}[chapter]}
\floatname{codelisting}{Listing}
\newcommand*\listoflistings{\listof{codelisting}{List of Listings}}
\makeatother
\makeatletter
\makeatother
\makeatletter
\@ifpackageloaded{caption}{}{\usepackage{caption}}
\@ifpackageloaded{subcaption}{}{\usepackage{subcaption}}
\makeatother

\ifLuaTeX
  \usepackage{selnolig}  % disable illegal ligatures
\fi
\usepackage{bookmark}

\IfFileExists{xurl.sty}{\usepackage{xurl}}{} % add URL line breaks if available
\urlstyle{same} % disable monospaced font for URLs
\hypersetup{
  pdftitle={PS 1},
  pdfauthor={Ben Schiffman},
  colorlinks=true,
  linkcolor={blue},
  filecolor={Maroon},
  citecolor={Blue},
  urlcolor={Blue},
  pdfcreator={LaTeX via pandoc}}


\title{PS 1}
\author{Ben Schiffman}
\date{2001-10-05}

\begin{document}
\maketitle


\begin{enumerate}
\def\labelenumi{\arabic{enumi}.}
\tightlist
\item
  \textbf{PS1:} Due Sat Oct 5 at 5:00PM Central. Worth 50 points.
  Initiate your
  \textbf{\href{https://classroom.github.com/a/uhUVze3Y}{repo}}
\end{enumerate}

We use (\texttt{*}) to indicate a problem that we think might be time
consuming.

Steps to submit (5 points on PS1 and 10 points on PS2)

\begin{enumerate}
\def\labelenumi{\arabic{enumi}.}
\tightlist
\item
  ``This submission is my work alone and complies with the 30538
  integrity policy.'' Add your initials to indicate your agreement:
  **\_\_**
\item
  ``I have uploaded the names of anyone I worked with on the problem set
  \textbf{\href{https://docs.google.com/forms/d/1-zzHx762odGlpVWtgdIC55vqF-j3gqdAp6Pno1rIGK0/edit}{here}}''
  **\_\_** (1 point)
\item
  Late coins used this pset: **\_\_** Late coins left after submission:
  **\_\_**
\item
  Knit your \texttt{ps1.qmd} to make \texttt{ps1.pdf}.

  \begin{itemize}
  \tightlist
  \item
    The PDF should not be more than 25 pages. Use \texttt{head()} and
    re-size figures when appropriate.
  \end{itemize}
\item
  Push \texttt{ps1.qmd} and \texttt{ps1.pdf} to your github repo. It is
  fine to use Github Desktop.
\item
  Submit \texttt{ps1.pdf} via Gradescope (4 points)
\item
  Tag your submission in Gradescope
\end{enumerate}

\subsection{Read in one percent sample (15
Points)}\label{read-in-one-percent-sample-15-points}

\begin{enumerate}
\def\labelenumi{\arabic{enumi}.}
\tightlist
\item
\end{enumerate}

\begin{Shaded}
\begin{Highlighting}[]
\ImportTok{import}\NormalTok{ pandas }\ImportTok{as}\NormalTok{ pd }
\ImportTok{import}\NormalTok{ os}
\ImportTok{import}\NormalTok{ time}


\NormalTok{pset\_dpath }\OperatorTok{=} \VerbatimStringTok{r"/Users/benschiffman/Desktop/Python 2/ppha30538\_fall2024/problem\_sets/ps1/data/"}

\NormalTok{oneperc\_path }\OperatorTok{=} \VerbatimStringTok{r"parking\_tickets\_one\_percent.csv"}

\NormalTok{start\_time }\OperatorTok{=}\NormalTok{ time.time()}
\NormalTok{df }\OperatorTok{=}\NormalTok{ pd.read\_csv(os.path.join(pset\_dpath, oneperc\_path))}
\NormalTok{end\_time }\OperatorTok{=}\NormalTok{ time.time()}
\NormalTok{total\_time }\OperatorTok{=}\NormalTok{ end\_time }\OperatorTok{{-}}\NormalTok{ start\_time}

\BuiltInTok{print}\NormalTok{(}\StringTok{"it takes"}\NormalTok{, total\_time, }\StringTok{"seconds to read this file"}\NormalTok{)}
\BuiltInTok{print}\NormalTok{(}\BuiltInTok{len}\NormalTok{(df))}
\ControlFlowTok{assert} \BuiltInTok{len}\NormalTok{(df) }\OperatorTok{==} \DecValTok{287458}\NormalTok{, }\StringTok{"Well, try again"}
\end{Highlighting}
\end{Shaded}

\begin{verbatim}
it takes 0.5482528209686279 seconds to read this file
287458
\end{verbatim}

\begin{verbatim}
/var/folders/yb/wxm346513_x57339c0hyc5280000gn/T/ipykernel_27905/3155636270.py:11: DtypeWarning: Columns (7) have mixed types. Specify dtype option on import or set low_memory=False.
  df = pd.read_csv(os.path.join(pset_dpath, oneperc_path))
\end{verbatim}

it takes 0.581108808517456 seconds to read this file. 1.

\begin{Shaded}
\begin{Highlighting}[]
\NormalTok{size }\OperatorTok{=}\NormalTok{ os.path.getsize(os.path.join(pset\_dpath, oneperc\_path))}
\BuiltInTok{print}\NormalTok{(size}\OperatorTok{/}\DecValTok{1000000}\NormalTok{)}
\end{Highlighting}
\end{Shaded}

\begin{verbatim}
83.655348
\end{verbatim}

The data dataset is 83.655348 MB 1.

\begin{Shaded}
\begin{Highlighting}[]
\NormalTok{df.head()}
\NormalTok{df.tail()}

\NormalTok{top\_data }\OperatorTok{=}\NormalTok{ df[:}\DecValTok{500}\NormalTok{]}

\NormalTok{top\_data[}\StringTok{"issue\_date"}\NormalTok{] }\OperatorTok{=}\NormalTok{ pd.to\_datetime(top\_data[}\StringTok{"issue\_date"}\NormalTok{])}

\CommentTok{\#this function compares dates in a column to see if they are in chronological order by index. }
\KeywordTok{def}\NormalTok{ is\_sorted(data):}
  \BuiltInTok{sorted} \OperatorTok{=} \VariableTok{False}
  \ControlFlowTok{for}\NormalTok{ date }\KeywordTok{in}\NormalTok{ data:}
\NormalTok{    baseline }\OperatorTok{=}\NormalTok{ data[}\DecValTok{0}\NormalTok{]}
    \ControlFlowTok{if}\NormalTok{ date }\OperatorTok{\textgreater{}=}\NormalTok{ baseline:}
\NormalTok{      baseline }\OperatorTok{=}\NormalTok{ date}
    \ControlFlowTok{else}\NormalTok{: }
      \ControlFlowTok{return} \BuiltInTok{print}\NormalTok{(}\StringTok{"The column is not sorted"}\NormalTok{)}
  \ControlFlowTok{return} \BuiltInTok{print}\NormalTok{(}\StringTok{"The column is sorted"}\NormalTok{)}
      
\NormalTok{is\_sorted(top\_data[}\StringTok{"issue\_date"}\NormalTok{])}
\end{Highlighting}
\end{Shaded}

\begin{verbatim}
The column is sorted
\end{verbatim}

\begin{verbatim}
/var/folders/yb/wxm346513_x57339c0hyc5280000gn/T/ipykernel_27905/2682147371.py:6: SettingWithCopyWarning: 
A value is trying to be set on a copy of a slice from a DataFrame.
Try using .loc[row_indexer,col_indexer] = value instead

See the caveats in the documentation: https://pandas.pydata.org/pandas-docs/stable/user_guide/indexing.html#returning-a-view-versus-a-copy
  top_data["issue_date"] = pd.to_datetime(top_data["issue_date"])
\end{verbatim}

It seems like the data is sorted by issue date and time. And the column
is sorted.

\subsection{Cleaning the data and benchmarking (15
Points)}\label{cleaning-the-data-and-benchmarking-15-points}

\begin{enumerate}
\def\labelenumi{\arabic{enumi}.}
\tightlist
\item
\end{enumerate}

\begin{Shaded}
\begin{Highlighting}[]
\NormalTok{df[}\StringTok{"issue\_date"}\NormalTok{] }\OperatorTok{=}\NormalTok{ pd.to\_datetime(df[}\StringTok{"issue\_date"}\NormalTok{])}
\NormalTok{df\_2017 }\OperatorTok{=}\NormalTok{ df[df[}\StringTok{"issue\_date"}\NormalTok{].dt.year }\OperatorTok{==} \DecValTok{2017}\NormalTok{]}
\NormalTok{ct }\OperatorTok{=} \BuiltInTok{len}\NormalTok{(df\_2017)}

\NormalTok{year\_ct }\OperatorTok{=}\NormalTok{ ct }\OperatorTok{*} \DecValTok{100}
\NormalTok{year\_ct}
\end{Highlighting}
\end{Shaded}

\begin{verbatim}
2236400
\end{verbatim}

There are 22364 tickets in the 1 percent dataset, which implies there
were \textasciitilde2,236,400 tickets issued in 2017. The article states
that illinois issues more than 3 million tickets a year. This indicates
that 2017 is a lower ticket year, but we need a lot more information to
make much more of a statement than that.

\begin{enumerate}
\def\labelenumi{\arabic{enumi}.}
\tightlist
\item
\end{enumerate}

\begin{Shaded}
\begin{Highlighting}[]
\ImportTok{import}\NormalTok{ altair }\ImportTok{as}\NormalTok{ alt}

\NormalTok{top\_violations }\OperatorTok{=}\NormalTok{ df.groupby(}\StringTok{"violation\_description"}\NormalTok{).size()}
\NormalTok{top\_violations }\OperatorTok{=}\NormalTok{ top\_violations.sort\_values(ascending}\OperatorTok{=}\VariableTok{False}\NormalTok{).head(}\DecValTok{20}\NormalTok{)}
\NormalTok{top\_violations }\OperatorTok{=}\NormalTok{ top\_violations.reset\_index()}
\NormalTok{top\_violations.columns }\OperatorTok{=}\NormalTok{ [}\StringTok{"violation\_description"}\NormalTok{, }\StringTok{"count"}\NormalTok{]}

\NormalTok{chart }\OperatorTok{=}\NormalTok{ alt.Chart(top\_violations).mark\_bar().encode(}
\NormalTok{    alt.X(}\StringTok{"violation\_description"}\NormalTok{, axis}\OperatorTok{=}\NormalTok{alt.Axis(labelAngle}\OperatorTok{={-}}\DecValTok{45}\NormalTok{, labelLimit }\OperatorTok{=} \DecValTok{400}\NormalTok{)),}
\NormalTok{    alt.Y(}\StringTok{"count"}\NormalTok{)}
\NormalTok{).properties(}
\NormalTok{    width}\OperatorTok{=}\DecValTok{500}\NormalTok{,}
\NormalTok{    title }\OperatorTok{=} \StringTok{"Count of Top 20 Violation Types"}
\NormalTok{)}
\NormalTok{chart}

\CommentTok{\#ChatGPT was used in this cell for ideation on how to sort data into top 20 violation types, as well as altair functionality specifically formatting labels. I used several queries, including "Can you provide me a list of altair tools that might be used to label formatting" after which it provided a list as well as example code, and "What I want to do is just pick the top 20 of this data type in the dataframe and then make the graph. How can I go about that" in which is provided several example paths and I chose to use groupby because I was already familiar with its functionality.}
\end{Highlighting}
\end{Shaded}

\begin{verbatim}
alt.Chart(...)
\end{verbatim}

\subsection{Visual Encoding (15
Points)}\label{visual-encoding-15-points}

\begin{enumerate}
\def\labelenumi{\arabic{enumi}.}
\tightlist
\item
\end{enumerate}

\begin{Shaded}
\begin{Highlighting}[]
\NormalTok{pd.set\_option(}\StringTok{\textquotesingle{}display.max\_columns\textquotesingle{}}\NormalTok{, }\VariableTok{None}\NormalTok{)}
\NormalTok{df.head()}
\end{Highlighting}
\end{Shaded}

\begin{longtable}[]{@{}lllllllllllllllllllllllll@{}}
\toprule\noalign{}
& Unnamed: 0 & ticket\_number & issue\_date & violation\_location &
license\_plate\_number & license\_plate\_state & license\_plate\_type &
zipcode & violation\_code & violation\_description & unit &
unit\_description & vehicle\_make & fine\_level1\_amount &
fine\_level2\_amount & current\_amount\_due & total\_payments &
ticket\_queue & ticket\_queue\_date & notice\_level &
hearing\_disposition & notice\_number & officer & address \\
\midrule\noalign{}
\endhead
\bottomrule\noalign{}
\endlastfoot
0 & 1 & 51482901.0 & 2007-01-01 01:25:00 & 5762 N AVONDALE &
d41ee9a4cb0676e641399ad14aaa20d06f2c6896de6366... & IL & PAS &
606184118.0 & 0964090E & RESIDENTIAL PERMIT PARKING & 16.0 & CPD & LEXU
& 50 & 100 & 0.0 & 50.0 & Paid & 2007-03-20 & DETR & Liable &
5.080059e+09 & 17266 & 5700 n avondale, chicago, il \\
1 & 2 & 50681501.0 & 2007-01-01 01:51:00 & 2724 W FARRAGUT &
3395fd3f71f18f9ea4f0a8e1f13bf0aa15052fc8e5605a... & IL & PAS &
606454911.0 & 0964090E & RESIDENTIAL PERMIT PARKING & 20.0 & CPD & FORD
& 50 & 100 & 0.0 & 50.0 & Paid & 2007-01-31 & VIOL & NaN & 5.079876e+09
& 10799 & 2700 w farragut, chicago, il \\
2 & 3 & 51579701.0 & 2007-01-01 02:22:00 & 1748 W ESTES &
302cb9c55f63ff828d7315c5589d97f1f8144904d66eb3... & IL & PAS &
604116803.0 & 0964150B & PARKING/STANDING PROHIBITED ANYTIME & 24.0 &
CPD & TOYT & 50 & 100 & 122.0 & 0.0 & Notice & 2007-02-28 & SEIZ & NaN &
5.037862e+09 & 17253 & 1700 w estes, chicago, il \\
3 & 4 & 51262201.0 & 2007-01-01 02:35:00 & 4756 N SHERIDAN &
94d018f52c7990cea326d1810a3278e2c6b1e8b44f3c52... & IL & PAS &
606601345.0 & 0976160F & EXPIRED PLATES OR TEMPORARY REGISTRATION & 23.0
& CPD & CHEV & 50 & 100 & 0.0 & 50.0 & Paid & 2007-01-11 & NaN & NaN &
5.075310e+09 & 3307 & 4700 n sheridan, chicago, il \\
4 & 5 & 51898001.0 & 2007-01-01 03:50:00 & 7134 S CAMPBELL &
876dd3a95179f4f1d720613f6e32a5a7b86b0e6f988bf4... & IL & PAS &
606291432.0 & 0964100A & WITHIN 15\textquotesingle{} OF FIRE HYDRANT &
8.0 & CPD & CHEV & 100 & 200 & 0.0 & 100.0 & Paid & 2007-04-25 & DETR &
NaN & 5.073568e+09 & 16820 & 7100 s campbell, chicago, il \\
\end{longtable}

\begin{longtable}[]{@{}ll@{}}
\toprule\noalign{}
Coloumn & Data Type \\
\midrule\noalign{}
\endhead
\bottomrule\noalign{}
\endlastfoot
Ticket Number & Ordinal \\
Issue Date & Temporal \\
Violation Location & Nominal \\
LP Number & Nominal \\
LP State & Nominal \\
LP Type & Nominal \\
ZIP Code & Nominal \\
Vio Code & Nominal \\
Vio Description & Nominal \\
Unit & Nominal \\
Unit Description & Nominal \\
Vehicle Make & Nominal \\
Fine Amount 1 & Quantitative \\
Fine Amount 2 & Quantitative \\
Amount Due & Quantitative \\
Total Payment & Quantitative \\
Ticket Queue & Ordinal \\
TQ Date & Temporal \\
Notice Level & Ordinal \\
Hearing Disp & Nominal \\
Notice Num & Nominal \\
Officer & Nominal \\
Address & Nominal \\
\end{longtable}

\begin{enumerate}
\def\labelenumi{\arabic{enumi}.}
\tightlist
\item
\end{enumerate}

\begin{Shaded}
\begin{Highlighting}[]
\NormalTok{df[}\StringTok{"Paid"}\NormalTok{] }\OperatorTok{=}\NormalTok{ df[}\StringTok{"ticket\_queue"}\NormalTok{].}\BuiltInTok{apply}\NormalTok{(}\KeywordTok{lambda}\NormalTok{ x: }\DecValTok{1} \ControlFlowTok{if}\NormalTok{ x }\OperatorTok{==} \StringTok{"Paid"} \ControlFlowTok{else} \DecValTok{0}\NormalTok{)}
\NormalTok{make\_paid }\OperatorTok{=}\NormalTok{ df.groupby([}\StringTok{"vehicle\_make"}\NormalTok{, }\StringTok{"Paid"}\NormalTok{]).size()}
\NormalTok{make\_paid }\OperatorTok{=}\NormalTok{ make\_paid.reset\_index(name }\OperatorTok{=} \StringTok{"count"}\NormalTok{)}
\NormalTok{pivot }\OperatorTok{=}\NormalTok{ make\_paid.pivot(index }\OperatorTok{=} \StringTok{"vehicle\_make"}\NormalTok{, columns }\OperatorTok{=} \StringTok{"Paid"}\NormalTok{, values }\OperatorTok{=} \StringTok{"count"}\NormalTok{).fillna(}\DecValTok{0}\NormalTok{)}
\NormalTok{pivot.columns }\OperatorTok{=}\NormalTok{ [}\StringTok{"Unpaid"}\NormalTok{, }\StringTok{"Paid"}\NormalTok{]}
\NormalTok{pivot.reset\_index(inplace}\OperatorTok{=}\VariableTok{True}\NormalTok{)}
\NormalTok{pivot}

\NormalTok{pivot[}\StringTok{"Ratio"}\NormalTok{] }\OperatorTok{=}\NormalTok{ pivot[}\StringTok{"Paid"}\NormalTok{]}\OperatorTok{/}\NormalTok{(pivot[}\StringTok{"Unpaid"}\NormalTok{] }\OperatorTok{+}\NormalTok{ pivot[}\StringTok{"Paid"}\NormalTok{]) }\OperatorTok{*} \DecValTok{100}
\NormalTok{pivot}

\NormalTok{paid\_chart }\OperatorTok{=}\NormalTok{ iol\_chart }\OperatorTok{=}\NormalTok{ alt.Chart(pivot).mark\_bar().encode(}
\NormalTok{    alt.X(}\StringTok{"vehicle\_make"}\NormalTok{),}
\NormalTok{    alt.Y(}\StringTok{"Ratio"}\NormalTok{)}
\NormalTok{).properties(}
\NormalTok{    title }\OperatorTok{=} \StringTok{"Percent of Tickets Paid by Make"}
\NormalTok{)}
\NormalTok{paid\_chart}

\CommentTok{\#chatGPT utilized for pivot functionality on this problem. I used a number of queries to ascertain how to use this pandas functionality to convert from the groupby object into a wide dataframe I could use for making the graph easily.}
\end{Highlighting}
\end{Shaded}

\begin{verbatim}
alt.Chart(...)
\end{verbatim}

\begin{enumerate}
\def\labelenumi{\arabic{enumi}.}
\tightlist
\item
\end{enumerate}

\begin{Shaded}
\begin{Highlighting}[]
\ImportTok{import}\NormalTok{ altair }\ImportTok{as}\NormalTok{ alt}
\ImportTok{import}\NormalTok{ datetime}

\NormalTok{df}
\NormalTok{df[}\StringTok{"issue\_date"}\NormalTok{] }\OperatorTok{=}\NormalTok{ pd.to\_datetime(df[}\StringTok{"issue\_date"}\NormalTok{])}
\NormalTok{df[}\StringTok{"Year"}\NormalTok{] }\OperatorTok{=}\NormalTok{ df[}\StringTok{"issue\_date"}\NormalTok{].dt.year}
\NormalTok{year\_tix }\OperatorTok{=}\NormalTok{ df.groupby(}\StringTok{"Year"}\NormalTok{).size().reset\_index(name }\OperatorTok{=} \StringTok{"Count"}\NormalTok{)}
\NormalTok{year\_tix.head()}

\NormalTok{alt.Chart(year\_tix).mark\_area(}
\NormalTok{    color}\OperatorTok{=}\StringTok{"lightblue"}\NormalTok{,}
\NormalTok{    interpolate}\OperatorTok{=}\StringTok{\textquotesingle{}step{-}after\textquotesingle{}}\NormalTok{,}
\NormalTok{    line}\OperatorTok{=}\VariableTok{True}
\NormalTok{).encode(}
\NormalTok{    x}\OperatorTok{=}\StringTok{\textquotesingle{}Year\textquotesingle{}}\NormalTok{,}
\NormalTok{    y}\OperatorTok{=}\StringTok{\textquotesingle{}Count\textquotesingle{}}
\NormalTok{)}
\end{Highlighting}
\end{Shaded}

\begin{verbatim}
alt.Chart(...)
\end{verbatim}

\begin{enumerate}
\def\labelenumi{\arabic{enumi}.}
\tightlist
\item
\end{enumerate}

\begin{Shaded}
\begin{Highlighting}[]
\ImportTok{import}\NormalTok{ altair }\ImportTok{as}\NormalTok{ alt}

\NormalTok{df}
\NormalTok{ticket\_counts }\OperatorTok{=}\NormalTok{ df.groupby(df[}\StringTok{"issue\_date"}\NormalTok{].dt.date).size().reset\_index(name}\OperatorTok{=}\StringTok{"count"}\NormalTok{)}
\NormalTok{ticket\_counts[}\StringTok{"issue\_date"}\NormalTok{] }\OperatorTok{=}\NormalTok{ pd.to\_datetime(ticket\_counts[}\StringTok{"issue\_date"}\NormalTok{])}

\NormalTok{alt.Chart(ticket\_counts, title }\OperatorTok{=} \StringTok{"Daily Ticket Counts"}\NormalTok{).mark\_rect().encode(}
\NormalTok{    alt.X(}\StringTok{"date(issue\_date):O"}\NormalTok{).title(}\StringTok{"Day"}\NormalTok{).axis(}\BuiltInTok{format}\OperatorTok{=}\StringTok{"}\SpecialCharTok{\%e}\StringTok{"}\NormalTok{, labelAngle}\OperatorTok{=}\DecValTok{0}\NormalTok{),}
\NormalTok{    alt.Y(}\StringTok{"month(issue\_date):O"}\NormalTok{).title(}\StringTok{"Month"}\NormalTok{),}
\NormalTok{    alt.Color(}\StringTok{"count"}\NormalTok{).title(}\VariableTok{None}\NormalTok{),}
\NormalTok{    tooltip}\OperatorTok{=}\NormalTok{[}
\NormalTok{        alt.Tooltip(}\StringTok{"monthdate(issue\_date)"}\NormalTok{, title}\OperatorTok{=}\StringTok{"Date"}\NormalTok{),}
\NormalTok{        alt.Tooltip(}\StringTok{"count"}\NormalTok{, title}\OperatorTok{=}\StringTok{"Tickets Count"}\NormalTok{),}
\NormalTok{    ],}
\NormalTok{).configure\_view(}
\NormalTok{    step}\OperatorTok{=}\DecValTok{13}\NormalTok{,}
\NormalTok{    strokeWidth}\OperatorTok{=}\DecValTok{0}
\NormalTok{).configure\_axis(}
\NormalTok{    domain}\OperatorTok{=}\VariableTok{False}
\NormalTok{)}

\CommentTok{"""}
\CommentTok{I used chatGPT for error handling, specifically for the error "TypeError: Object of type date is not JSON serializable." ChatGPT suggested that the "issue\_date" column may not be a datetime object and needed conversion. }
\CommentTok{"""}
\end{Highlighting}
\end{Shaded}

\begin{verbatim}
'\nI used chatGPT for error handling, specifically for the error "TypeError: Object of type date is not JSON serializable." ChatGPT suggested that the "issue_date" column may not be a datetime object and needed conversion. \n'
\end{verbatim}

\begin{enumerate}
\def\labelenumi{\arabic{enumi}.}
\tightlist
\item
\end{enumerate}

\begin{Shaded}
\begin{Highlighting}[]
\ImportTok{import}\NormalTok{ altair }\ImportTok{as}\NormalTok{ alt}
\ImportTok{import}\NormalTok{ datetime }\ImportTok{as}\NormalTok{ dt}

\CommentTok{\#first identify top 5 violation types}
\NormalTok{top5\_violations }\OperatorTok{=}\NormalTok{ df.copy()}
\NormalTok{top5\_violations }\OperatorTok{=}\NormalTok{ top5\_violations.groupbytop5\_violations }\OperatorTok{=}\NormalTok{ top5\_violations.groupby(}\StringTok{"violation\_description"}\NormalTok{).size().reset\_index(name }\OperatorTok{=} \StringTok{"count"}\NormalTok{)}
\NormalTok{top5\_violations }\OperatorTok{=}\NormalTok{ top5\_violations.sort\_values(by }\OperatorTok{=} \StringTok{"count"}\NormalTok{, ascending}\OperatorTok{=}\VariableTok{False}\NormalTok{).head(}\DecValTok{5}\NormalTok{)}
\NormalTok{top5\_violations}

\CommentTok{\#create a new df with just top 5}
\NormalTok{lasagna\_df }\OperatorTok{=}\NormalTok{ pd.merge(df, top5\_violations, on }\OperatorTok{=} \StringTok{"violation\_description"}\NormalTok{, how }\OperatorTok{=} \StringTok{"right"}\NormalTok{)}
\NormalTok{lasagna\_df[}\StringTok{"issue\_date"}\NormalTok{] }\OperatorTok{=}\NormalTok{ pd.to\_datetime(lasagna\_df[}\StringTok{"issue\_date"}\NormalTok{])}
\NormalTok{lasagna\_df[}\StringTok{\textquotesingle{}year\_month\textquotesingle{}}\NormalTok{] }\OperatorTok{=}\NormalTok{ lasagna\_df[}\StringTok{\textquotesingle{}issue\_date\textquotesingle{}}\NormalTok{].dt.to\_period(}\StringTok{\textquotesingle{}M\textquotesingle{}}\NormalTok{) }
\NormalTok{lasagna\_df}

\CommentTok{\#summarize data becuase of too many rows}
\CommentTok{\#ChatGPT aided with this step. I used the prompt "How to deal with max rows error altair" and further explored the option of aggregation with the model}

\NormalTok{monthly\_counts }\OperatorTok{=}\NormalTok{ lasagna\_df.groupby([}\StringTok{"violation\_description"}\NormalTok{,}\StringTok{\textquotesingle{}year\_month\textquotesingle{}}\NormalTok{]).size().reset\_index(name}\OperatorTok{=}\StringTok{\textquotesingle{}count\textquotesingle{}}\NormalTok{)}
\NormalTok{monthly\_counts[}\StringTok{"year\_month"}\NormalTok{] }\OperatorTok{=}\NormalTok{ monthly\_counts[}\StringTok{"year\_month"}\NormalTok{].dt.to\_timestamp()}
\NormalTok{monthly\_counts}

\CommentTok{\#create graph}
\NormalTok{color\_condition }\OperatorTok{=}\NormalTok{ alt.condition(}
    \StringTok{"month(datum.value) == 1 \&\& date(datum.value) == 1"}\NormalTok{,}
\NormalTok{    alt.value(}\StringTok{"black"}\NormalTok{),}
\NormalTok{    alt.value(}\VariableTok{None}\NormalTok{),}
\NormalTok{)}
\NormalTok{alt.Chart(monthly\_counts, width}\OperatorTok{=}\DecValTok{300}\NormalTok{, height}\OperatorTok{=}\DecValTok{100}\NormalTok{).transform\_filter(}
\NormalTok{    alt.datum.symbol }\OperatorTok{!=} \StringTok{"GOOG"}
\NormalTok{).mark\_rect().encode(}
\NormalTok{    alt.X(}\StringTok{"yearmonthdate(year\_month):O"}\NormalTok{)}
\NormalTok{        .title(}\StringTok{"Time"}\NormalTok{)}
\NormalTok{        .axis(}
            \BuiltInTok{format}\OperatorTok{=}\StringTok{"\%Y"}\NormalTok{,}
\NormalTok{            labelAngle}\OperatorTok{=}\DecValTok{0}\NormalTok{,}
\NormalTok{            labelOverlap}\OperatorTok{=}\VariableTok{False}\NormalTok{,}
\NormalTok{            labelColor}\OperatorTok{=}\NormalTok{color\_condition,}
\NormalTok{            tickColor}\OperatorTok{=}\NormalTok{color\_condition,}
\NormalTok{        ),}
\NormalTok{    alt.Y(}\StringTok{"violation\_description:N"}\NormalTok{).title(}\VariableTok{None}\NormalTok{),}
\NormalTok{    alt.Color(}\StringTok{"sum(count)"}\NormalTok{).title(}\StringTok{"Tickets Issued"}\NormalTok{)}
\NormalTok{)}
\end{Highlighting}
\end{Shaded}

\begin{verbatim}
/var/folders/yb/wxm346513_x57339c0hyc5280000gn/T/ipykernel_27905/383925509.py:6: UserWarning: Pandas doesn't allow columns to be created via a new attribute name - see https://pandas.pydata.org/pandas-docs/stable/indexing.html#attribute-access
  top5_violations = top5_violations.groupbytop5_violations = top5_violations.groupby("violation_description").size().reset_index(name = "count")
\end{verbatim}

\begin{verbatim}
alt.Chart(...)
\end{verbatim}

\begin{Shaded}
\begin{Highlighting}[]
\ImportTok{from}\NormalTok{ vega\_datasets }\ImportTok{import}\NormalTok{ data}

\NormalTok{source }\OperatorTok{=}\NormalTok{ data.stocks()}
\NormalTok{source}
\NormalTok{color\_condition }\OperatorTok{=}\NormalTok{ alt.condition(}
    \StringTok{"month(datum.value) == 1 \&\& date(datum.value) == 1"}\NormalTok{,}
\NormalTok{    alt.value(}\StringTok{"black"}\NormalTok{),}
\NormalTok{    alt.value(}\VariableTok{None}\NormalTok{),}
\NormalTok{)}

\NormalTok{alt.Chart(source, width}\OperatorTok{=}\DecValTok{300}\NormalTok{, height}\OperatorTok{=}\DecValTok{100}\NormalTok{).transform\_filter(}
\NormalTok{    alt.datum.symbol }\OperatorTok{!=} \StringTok{"GOOG"}
\NormalTok{).mark\_rect().encode(}
\NormalTok{    alt.X(}\StringTok{"yearmonthdate(date):O"}\NormalTok{)}
\NormalTok{        .title(}\StringTok{"Time"}\NormalTok{)}
\NormalTok{        .axis(}
            \BuiltInTok{format}\OperatorTok{=}\StringTok{"\%Y"}\NormalTok{,}
\NormalTok{            labelAngle}\OperatorTok{=}\DecValTok{0}\NormalTok{,}
\NormalTok{            labelOverlap}\OperatorTok{=}\VariableTok{False}\NormalTok{,}
\NormalTok{            labelColor}\OperatorTok{=}\NormalTok{color\_condition,}
\NormalTok{            tickColor}\OperatorTok{=}\NormalTok{color\_condition,}
\NormalTok{        ),}
\NormalTok{    alt.Y(}\StringTok{"symbol:N"}\NormalTok{).title(}\VariableTok{None}\NormalTok{),}
\NormalTok{    alt.Color(}\StringTok{"sum(price)"}\NormalTok{).title(}\StringTok{"Price"}\NormalTok{)}
\NormalTok{)}
\end{Highlighting}
\end{Shaded}

\begin{verbatim}
alt.Chart(...)
\end{verbatim}

\begin{enumerate}
\def\labelenumi{\arabic{enumi}.}
\tightlist
\item
  The prior three questions created the filled step, heat map, and
  lasagna charts. The filled step shows broader timeseries data, which
  is best for showing how one data point changes over time. The heat map
  chart is going to be best for comparing patterns over the course of
  weeks and days in a month as well as over the course of a year,
  aggreggated from many years of data. This is best for seeing a pattern
  within a calandar year for one data point, perhaps to establish a
  baseline pattern to compare an individual year to. Finally the lasagna
  chart compares patterns over time between multiple categories. This is
  best for comparing categories with eachother within the same time
  frame, and can show some basic patterns within each category.
\item
  To show a reader that distributions are not evenly distributed over
  time, I would show the heat map. This shows the patterns in
  distribution by day and month. That said, both other graphs show
  different time periods in different detail, so they have their own
  benefits. The heat map simply shows differences in distribution of all
  tickets over at the course of the year. The filled step graph shows a
  broad pattern, but may be due to external factors like policy changes,
  and the lasagna map may on the other hand show too small a sample of
  types of tickets to assume a wider pattern.
\end{enumerate}




\end{document}
